\BourbakiExercisesSection

\begin{enumerate}
\def\labelenumi{\arabic{enumi}.}
\item
  In a matrix U of order n, if, for each index i, the column of index i is replaced by the sum of the columns of index \(\neq i\) , the determinant of the matrix obtained is equal to \((-1)^{n-1}(n-1)\det U\) . If in U, for each i, the sum of the columns of index \(\neq i\) is subtracted from the column of index i, the determinant obtained is equal to \(-(n-2)2^{n-1}\det U\) .
\item
  Let \(\Delta = \det (\alpha_{ij})\) be the determinant of a matrix of order \(n\) ; for
\end{enumerate}

\[
1 \leqslant i \leqslant n - 1
\]

and \(1 \leqslant j \leqslant n - 1\) , we write

\[
\beta_ {i j} = \left| \begin{array}{c c} \alpha_ {\mathbf {1} j} & \alpha_ {\mathbf {1}, j + \mathbf {1}} \\ \alpha_ {i + \mathbf {1}, j} & \alpha_ {i + \mathbf {1}, j + \mathbf {1}} \end{array} \right|.
\]

Show that the determinant \(\det(\beta_{ij})\) of order n - 1 is equal to

\[
\alpha_ {1 2} \alpha_ {1 3} \dots \alpha_ {1, n - 1} \Delta .
\]

\begin{enumerate}
\def\labelenumi{\arabic{enumi}.}
\setcounter{enumi}{2}
\tightlist
\item
  Show the identity
\end{enumerate}

\[
\left| \begin{array}{cccccc} x _ {1} & y _ {1} & \alpha_ {1 3} & \alpha_ {1 4} & \ldots & \alpha_ {1 n} \\ \lambda_ {1} x _ {2} & x _ {2} & y _ {2} & \alpha_ {2 4} & \ldots & \alpha_ {2 n} \\ \lambda_ {1} \lambda_ {2} x _ {3} & \lambda_ {2} x _ {3} & x _ {3} & y _ {3} & \ldots & \alpha_ {3 n} \\ \cdot & \cdot & \cdot & \cdot & \cdot & \cdot \\ \lambda_ {1} \lambda_ {2} \ldots \lambda_ {n - 1} x _ {n} & \lambda_ {2} \ldots \lambda_ {n - 1} x _ {n} & \lambda_ {3} \ldots \lambda_ {n - 1} x _ {n} & \lambda_ {4} \ldots \lambda_ {n - 1} x _ {n} & \ldots & x _ {n} \\ \end{array} \right| = (x _ {1} - \lambda_ {1} y _ {1}) (x _ {2} - \lambda_ {2} y _ {2}) \ldots (x _ {n - 1} - \lambda_ {n - 1} y _ {n - 1}) x _ {n}.
\]

Deduce the following identities:

\[
\left| \begin{array}{c c c c c} 1 & 1 & 1 & \dots & 1 \\ b _ {1} & a _ {1} & a _ {1} & \dots & a _ {1} \\ b _ {1} & b _ {2} & a _ {2} & \dots & a _ {2} \\ \cdot & \cdot & \cdot & \cdot & \cdot \\ b _ {1} & b _ {2} & b _ {3} & \dots & a _ {n} \end{array} \right| = (a _ {1} - b _ {1}) (a _ {2} - b _ {2}) \dots (a _ {n} - b _ {n})
\]

\[
\left| \begin{array}{c c c c c} a _ {1} b _ {1} & a _ {1} b _ {2} & a _ {1} b _ {3} & \ldots & a _ {1} b _ {n} \\ a _ {1} b _ {2} & a _ {2} b _ {2} & a _ {2} b _ {3} & \ldots & a _ {2} b _ {n} \\ a _ {1} b _ {3} & a _ {2} b _ {3} & a _ {3} b _ {3} & \ldots & a _ {3} b _ {n} \\ \cdot \cdot & \cdot \cdot & \cdot \cdot & \cdot \cdot & \cdot \cdot \\ a _ {1} b _ {n} & a _ {2} b _ {n} & a _ {3} b _ {n} & \ldots & a _ {n} b _ {n} \end{array} \right| = a _ {1} b _ {n} (a _ {2} b _ {1} - a _ {1} b _ {2}) \ldots (a _ {n} b _ {n - 1} - a _ {n - 1} b _ {n})
\]

\[
\left| \begin{array}{cccccc} a _ {2} a _ {3} \ldots a _ {n} & a _ {3} a _ {4} \ldots a _ {n} b _ {1} & a _ {4} \ldots a _ {n} b _ {1} b _ {2} & \ldots & b _ {1} b _ {2} b _ {3} \ldots b _ {n - 1} \\ b _ {2} b _ {3} \ldots b _ {n} & a _ {3} a _ {4} \ldots a _ {n} a _ {1} & a _ {4} \ldots a _ {n} a _ {1} b _ {2} & \ldots & a _ {1} b _ {2} b _ {3} \ldots b _ {n - 1} \\ a _ {2} b _ {3} \ldots b _ {n} & b _ {3} b _ {4} \ldots b _ {n} b _ {1} & a _ {4} \ldots a _ {n} a _ {1} a _ {2} & \ldots & a _ {1} a _ {2} b _ {3} \ldots b _ {n - 1} \\ \cdot & \cdot & \cdot & \cdot & \cdot \\ a _ {2} a _ {3} \ldots b _ {n} & a _ {3} a _ {4} \ldots b _ {n} b _ {1} & a _ {4} \ldots b _ {n} b _ {1} b _ {2} & \ldots & a _ {1} a _ {2} a _ {3} \ldots a _ {n - 1} \\ \end{array} \right| = (a _ {1} a _ {2} \ldots a _ {n} - b _ {1} b _ {2} \ldots b _ {n}) ^ {n - 1}.
\]

\[
\left| \begin{array}{c c c c c c} x & a _ {1} & a _ {2} & \ldots & a _ {n - 1} & 1 \\ a _ {1} & x & a _ {2} & \ldots & a _ {n - 1} & 1 \\ a _ {1} & a _ {2} & x & \ldots & a _ {n - 1} & 1 \\ . & . & . & . & . & . \\ a _ {1} & a _ {2} & a _ {3} & \ldots & x & 1 \\ a _ {1} & a _ {2} & a _ {3} & \ldots & a _ {n} & 1 \end{array} \right| = (x - a _ {1}) (x - a _ {2}) \ldots (x - a _ {n})
\]

\[
\begin{array}{r l} & {\left| \begin{array}{l l l l l} x & a _ {1} & a _ {2} & \ldots & a _ {n} \\ a _ {1} & x & a _ {2} & \ldots & a _ {n} \\ a _ {1} & a _ {2} & x & \ldots & a _ {n} \\ \cdot & \cdot & \cdot & \cdot & \cdot \\ a _ {1} & a _ {2} & a _ {3} & \ldots & x \end{array} \right|} \\ & {\qquad = (x + a _ {1} + a _ {2} + \dots + a _ {n}) (x - a _ {1}) (x - a _ {2}) \ldots (x - a _ {n})} \end{array}
\]

(reduce the last determinant to the preceding one).

\begin{enumerate}
\def\labelenumi{\arabic{enumi}.}
\setcounter{enumi}{3}
\tightlist
\item
  Calculate the determinant
\end{enumerate}

\[
\Delta_ {n} = \left| \begin{array}{c c c c c} a _ {1} + b _ {1} & b _ {1} & b _ {1} & \ldots & b _ {1} \\ b _ {2} & a _ {2} + b _ {2} & b _ {2} & \ldots & b _ {2} \\ \cdot & \cdot & \cdot & \cdot & \cdot \\ b _ {n} & b _ {n} & b _ {n} & \ldots & a _ {n} + b _ {n} \end{array} \right|
\]

(express \(\Delta_{n}\) in terms of \(\Delta_{n-1}\) ).

\begin{enumerate}
\def\labelenumi{\arabic{enumi}.}
\setcounter{enumi}{4}
\tightlist
\item
  If the \(a_i\) and \(b_j\) are elements of a commutative field such that \(a_i + b_j \neq 0\) for every ordered pair of indices \((i,j)\) , prove that
\end{enumerate}

\[
\det \left(\frac {1}{a _ {i} + b _ {j}}\right) = \frac {\sum_ {i <   j} (a _ {j} - a _ {i}) (b _ {j} - b _ {i})}{\sum_ {i , j} (a _ {i} + b _ {j})}
\]

(``Cauchy's determinant'').

\begin{enumerate}
\def\labelenumi{\arabic{enumi}.}
\setcounter{enumi}{5}
\tightlist
\item
  Show that, if X is a matrix with n rows and m columns, Y a matrix with p rows and n columns and Z = YX the product matrix with p rows and m columns, the minors of Z of order q are zero if n \textless{} q; if \(q \leqslant n\) , they are given by
\end{enumerate}

\[
\det (Z _ {L, H}) = \sum_ {K} \det (Y _ {L, K}) \det (X _ {K, H})
\]

where K runs through the set of subsets of \([1, n]\) with q elements (use formula (3) of §7, no. 2).

\begin{enumerate}
\def\labelenumi{\arabic{enumi}.}
\setcounter{enumi}{6}
\item
  Let \(\Delta = \det(\alpha_{ij})\) be the determinant of a matrix U of order n; for each index i ( \(1 \leqslant i \leqslant n\) ) let \(\Delta_i\) denote the determinant obtained by multiplying the element \(\alpha_{ij}\) in U by \(\beta_j\) for \(1 \leqslant j \leqslant n\) . Show that the sum of the n determinants \(\Delta_i\) ( \(1 \leqslant i \leqslant n\) ) is equal to \((\beta_1 + \beta_2 + \cdots + \beta_n)\Delta\) (expand \(\Delta_i\) by the row of index i).
\item
  Let \(\Delta = \det(\alpha_{ij})\) be the determinant of a matrix \(U\) of order \(n\) and \(\sigma\) a permutation belonging to \(\mathfrak{S}_n\) ; for each index \(i\) ( \(1 \leqslant i \leqslant n\) ) let \(\Delta_i\) be the determinant obtained by replacing the element \(\alpha_{ij}\) in \(U\) by \(\alpha_{i,\sigma(j)}\) for \(1 \leqslant j \leqslant n\) ; if \(p\) is the number of indices invariant under the permutation \(\sigma\) , show that the sum of the \(n\) determinants \(\Delta_i\) ( \(1 \leqslant i \leqslant n\) ) is equal to \(p\Delta\) (same method as in Exercise 7).
\item
  Let A be a square matrix of order n, B a submatrix of A with p rows and q columns and C the matrix obtained by multiplying in A, each of the elements belonging to B by the same scalar \(\alpha\) . Show that each term in the total expansion of det C is equal to the corresponding term in the total expansion of det A multiplied by a scalar of the form \(\alpha^{r}\) , where \(r \geqslant p + q - n\) and r depends on the term considered (perform a suitable Laplace expansion of det C).
\item
  Let \(\Gamma, \Delta\) be the determinants of two matrices \(U\) , \(V\) of order \(n\) ; let H and K be any two subsets of \([1, n]\) with \(p\) elements and \((i_k)\) (resp. \((j_k)\) ) the sequence obtained by arranging the elements of H (resp. K) in increasing order; let \(\Gamma_{H,K}\) be the determinant obtained by replacing in \(U\) the column of index \(i_k\) by the column of \(V\) of index \(j_k\) , for \(1 \leqslant k \leqslant p\) ; similarly let \(\Delta_{K,H}\) be the determinant obtained by replacing in \(V\) the column of index \(j_k\) by the column of \(U\) of index \(i_k\) , for \(1 \leqslant k \leqslant p\) . Show that, for every subset \(H\) of \([1, n]\) with \(p\) elements,
\end{enumerate}

\[
\Gamma \Delta = \sum_ {\mathbf {K}} \Gamma_ {\mathbf {H}, \mathbf {K}} \Delta_ {\mathbf {K}, \mathbf {H}}
\]

where K runs through the set of subsets of \([1, n]\) with p elements (use formulae (21) and (22) of no. 6).

\begin{enumerate}
\def\labelenumi{\arabic{enumi}.}
\setcounter{enumi}{10}
\tightlist
\item
  Let \(\Delta\) be the determinant of a square matrix \(X\) of order \(n\) and \(\Delta_p\) the determinant of the square matrix \(\bigwedge^p X\) , of order \(\binom{n}{p}\) . Show that
\end{enumerate}

\[
\Delta_ {p} \Delta_ {n - p} = \Delta_ {(p)} ^ {(n)}
\]

(use formulae (21) and (22) of no. 6).

\begin{enumerate}
\def\labelenumi{\arabic{enumi}.}
\setcounter{enumi}{11}
\tightlist
\item
  A matrix \((\alpha_{ij})\) of order \(n\) is called centrosymmetric (resp. skew-centrosymmetric) if \(\alpha_{n - i + 1, n - j + 1} = \alpha_{ij}\) (resp. \(\alpha_{n - i + 1, n - j + 1} = -\alpha_{ij}\) ) for \(1 \leqslant i \leqslant n\) , \(1 \leqslant j \leqslant n\) .
\end{enumerate}

\begin{enumerate}
\def\labelenumi{(\alph{enumi})}
\item
  Show that the determinant of a centrosymmetric matrix of even order 2p can be expressed in the form of a product of two determinants of order p and the determinant of a centrosymmetric matrix of order \(2p + 1\) in the form of a product of a determinant of order p and a determinant of order \(p + 1\) .
\item
  Show that the determinant of a skew-centrosymmetric matrix of even order \(2p\) can be expressed in the form of a product of two determinants of order \(p\) . The determinant of a skew-centrosymmetric matrix of odd order \(2p + 1\) is zero if in A the relation \(2\xi = 0\) implies \(\xi = 0\) ; otherwise it can be expressed in the form of the product of \(\alpha_{p+1,p+1}\) by two determinants of order \(p\) .
\end{enumerate}

\begin{enumerate}
\def\labelenumi{\arabic{enumi}.}
\setcounter{enumi}{12}
\tightlist
\item
  Let \(\Delta = \det(\alpha_{ij})\) be a determinant of order \(n\) and \(\Delta_{ij}\) the minor of order \(n - 1\) , the complement of \(\alpha_{ij}\) . Show that
\end{enumerate}

\[
\left| \begin{array}{c c c c c} \alpha_ {1 1} & \alpha_ {1 2} & \ldots & \alpha_ {1 n} & x _ {1} \\ \alpha_ {2 1} & \alpha_ {2 2} & \ldots & \alpha_ {2 n} & x _ {2} \\ \cdot & \cdot & \cdot & \cdot & \cdot \\ \alpha_ {n 1} & \alpha_ {n 2} & \ldots & \alpha_ {n n} & x _ {n} \\ y _ {1} & y _ {2} & \ldots & y _ {n} & z \end{array} \right| = \Delta z - \sum_ {i, j} (- 1) ^ {i + j} \Delta_ {i j} x _ {i} y _ {j}.
\]

If \(\Delta = 0\) and the \(\alpha_{ij}\) belong to a field, show that the above determinant is the product of a linear form in \(x_1, x_2, \ldots, x_n\) and a linear form in \(y_1, y_2, \ldots, y_n\) (use Exercise 5 of §5 and Exercise 8 of II, §6).

Give an example where this result fails to hold when the ring of scalars A is not a field (take A to be the ring \(\mathbf{Z} / (6)\) and \(n = 2\) ).

\begin{enumerate}
\def\labelenumi{\arabic{enumi}.}
\setcounter{enumi}{13}
\tightlist
\item
  Show the identity
\end{enumerate}

\[
\left| \begin{array}{cccccc} 0 & 1 & 1 & 1 & \ldots & 1 \\ 1 & 0 & a _ {1} + a _ {2} & a _ {1} + a _ {3} & \ldots & a _ {1} + a _ {n} \\ 1 & a _ {2} + a _ {1} & 0 & a _ {2} + a _ {3} & \ldots & a _ {2} + a _ {n} \\ 1 & a _ {3} + a _ {1} & a _ {3} + a _ {2} & 0 & \ldots & a _ {3} + a _ {n} \\ \cdot & \cdot & \cdot & \cdot & \cdot & \cdot \\ 1 & a _ {n} + a _ {1} & a _ {n} + a _ {2} & a _ {n} + a _ {3} & \ldots & 0 \\ \end{array} \right| = (- 1) ^ {n} 2 ^ {n - 1} \sum_ {i = 1} ^ {n} a _ {1} a _ {2} \ldots a _ {i - 1} a _ {i + 1} \ldots a _ {n}
\]

(use Exercise 13).

\begin{enumerate}
\def\labelenumi{\arabic{enumi}.}
\setcounter{enumi}{14}
\item
  Show that if the columns of a square matrix of order n over A are linearly independent, the rows of the matrix are also linearly independent (cf.~II, § 10, Exercise 3).
\item
  Let E, F be two free A-modules of respective dimensions m and n.~For a linear mapping \(u: E \rightarrow F\) to be injective, it is necessary and sufficient that \(m \leqslant n\) and that, if X denotes the matrix of u with respect to any two bases of E and F, there exist no scalar \(\mu \neq 0\) such that the products by \(\mu\) of all the minors of X of order m are zero.
\end{enumerate}

\begin{enumerate}
\def\labelenumi{\arabic{enumi}.}
\setcounter{enumi}{16}
\tightlist
\item
  Let
\end{enumerate}

\[
\sum_ {j = 1} ^ {n} \alpha_ {i j} \xi_ {j} = \beta_ {i} \quad (1 \leqslant i \leqslant m)\tag{*}
\]

be a system of m linear equations in n unknowns over a commutative ring A. Let \(x_{j}\) ( \(1 \leqslant j \leqslant n\) ) be the columns of the matrix \(X = (\alpha_{ij})\) of this system and \(y = (\beta_{i})\) ; suppose that in X all the minors of order \textgreater p are zero but that \(x_{1} \wedge x_{2} \wedge \cdots \wedge x_{p} \neq 0\) . For the system (*) to have a solution, it is necessary that \(x_{1} \wedge x_{2} \wedge \cdots \wedge x_{p} \wedge y = 0\) . Conversely, if this condition holds, there exist \(n + 1\) elements \(\xi_{j}\) ( \(1 \leqslant j \leqslant n + 1\) ) of A such that \(\xi_{n+1} \neq 0\) and

\[
\sum_ {j = 1} ^ {n} \alpha_ {i j} \xi_ {j} = \beta_ {i} \xi_ {n + 1} \quad (1 \leqslant i \leqslant m).
\]

\begin{enumerate}
\def\labelenumi{\arabic{enumi}.}
\setcounter{enumi}{17}
\tightlist
\item
  Let E and F be two free A-modules of respective dimensions m and n, \(u:E \rightarrow F\) a linear mapping and X the matrix of u with respect to any two bases of E and F.
\end{enumerate}

\begin{enumerate}
\def\labelenumi{(\alph{enumi})}
\item
  If \(m \geqslant n\) and there exists a minor of \(X\) of order \(n\) which is invertible, \(u\) is surjective.
\item
  Conversely, if \(u\) is surjective, show that \(m \geqslant n\) and that there exists a minor of \(X\) of order \(n\) which is non-zero. If further, in the ring \(A\) , the ideal generated by the non-invertible elements is distinct from \(A\) , there exists a minor of \(X\) of order \(n\) which is invertible (consider the exterior power \(\bigwedge^{n}u\) ).
\item
  Let B be a commutative ring and A the product ring \(B \times B\) . Give an example of a surjective linear mapping of the A-module \(A^{2}\) onto A, all the elements of whose matrix (with respect to the canonical bases of \(A^{2}\) and A) are divisors of zero.
\end{enumerate}

\begin{enumerate}
\def\labelenumi{\arabic{enumi}.}
\setcounter{enumi}{18}
\item
  Let X be a matrix over a commutative field. For X to be of rank p, it suffices that there exist a minor of X of order p which is \(\neq0\) and such that all the minors of order \(p+1\) containing this minor of order p are zero (show that every column of X is then a linear combination of the p columns to which the elements of the minor of order p in question belong).
\item
  Let \(\mathbf{A}\) be a commutative ring, \(\mathbf{M}\) an arbitrary A-module and \(U = (\alpha_{ij})\) a matrix of type \((m, n)\) with elements in \(\mathbf{A}\) .
\end{enumerate}

\begin{enumerate}
\def\labelenumi{(\alph{enumi})}
\item
  Suppose that \(m = n\) and let \(\Delta = \det(U)\) ; then, if \(x_1, \ldots, x_n\) are elements of \(M\) such that \(\sum_{j=1}^{n} \alpha_{ij} x_j = 0\) for \(1 \leqslant i \leqslant n\) , then \(\Delta x_i = 0\) for every index \(i\) .
\item
  Suppose that \(m\) and \(n\) are arbitrary, that all the minors of \(U\) of order \(> r\) are zero and that a minor of order \(r\) is invertible, for example that of the matrix \((\alpha_{ij})\) where \(i \leqslant r\) and \(j \leqslant r\) ; then if \(a_i\) denotes the row of index \(i\) in \(U\) , the \(a_i\) of index \(\leqslant r\) form a basis of the submodule of \(A^n\) generated by all the rows of \(U\) ; for \(i > r\) , therefore \(a_i = \sum_{k=1}^{r} \lambda_{ik} a_k\) . Then let \(y_1, \ldots, y_m\) be elements of \(M\) ; for there to exist elements \(x_1, \ldots, x_n\) of \(M\) satisfying the system of equations \(\sum_{j=1}^{n} \alpha_{ij} x_j = y_i\) for \(1 \leqslant i \leqslant m\) , it is necessary and sufficient that \(y_i = \sum_{k=1}^{r} \lambda_{ik} y_k\) for \(i > r\) .
\end{enumerate}

¶ 21. Let K be an infinite commutative field, m, n, r three integers \(\geqslant0\) such that \(r\leqslant n\leqslant m\) , \(\mathbf{M}_{m,n}(\mathbf{K})\) the vector K-space (of dimension mn) consisting of the matrices of type \((m,n)\) over K and V a vector subspace of \(\mathbf{M}_{m,n}(\mathbf{K})\) such that r is the greatest value of the rank of matrices \(X\in V\) . We propose to prove the inequality

\[
\dim (\mathbf {V}) \leqslant m r.\tag{*}
\]

\begin{enumerate}
\def\labelenumi{(\alph{enumi})}
\tightlist
\item
  Show that we can assume that \(m = n\) and that the matrix
\end{enumerate}

\[
X _ {0} = \left( \begin{array}{c c} I _ {r} & 0 \\ 0 & 0 \end{array} \right)
\]

belongs to V. Deduce that every matrix \(X \in \mathbf{V}\) is then of the form

\[
\left( \begin{array}{c c} X _ {1 1} & X _ {1 2} \\ X _ {2 1} & 0 \end{array} \right)
\]

(where \(X_{11}\) is of order \(r\) ) with the condition \(X_{21}X_{12} = 0\) (use the fact that for all \(\xi \in \mathbf{K}\) , the matrix \(\xi X_0 + X\) is of rank \(\leqslant r\) ).

\begin{enumerate}
\def\labelenumi{(\alph{enumi})}
\setcounter{enumi}{1}
\tightlist
\item
  Deduce from (a) that, for two matrices \(X\) , \(Y\) in \(\mathbf{V}\) ,
\end{enumerate}

\[
X _ {2 1} Y _ {1 2} + Y _ {2 1} X _ {1 2} = 0.
\]

\begin{enumerate}
\def\labelenumi{(\alph{enumi})}
\setcounter{enumi}{2}
\tightlist
\item
  For every matrix \(X \in \mathbf{V}\) , we write \(F(X) = (X_{11}X_{12}) \in \mathbf{M}_{r,n}(\mathbf{K})\) ; \(f\) is a linear mapping and \(\dim(\mathbf{V}) = \dim(\operatorname{Im}(f)) + \dim(\operatorname{Ker}(f))\) . On the other hand, let \(u_x\) denote the linear form \((\mathrm{Y}_{11}\mathrm{Y}_{12}) \mapsto \operatorname{Tr}(\mathrm{X}_{21}\mathrm{Y}_{12})\) on \(\mathbf{M}_{r,n}(\mathbf{K})\) . Show that the linear mapping \(X \mapsto u_X\) of \(\operatorname{Ker}(f)\) into the dual \((\mathbf{M}_{r,n}(\mathbf{K}))^*\) is injective; complete the proof of (*) by noting that the image under \(X \mapsto u_X\) of \(\operatorname{Ker}(f)\) is contained in the orthogonal of \(\operatorname{Im}(f)\) .
\end{enumerate}

¶ 22. Let F be a mapping of \(\mathbf{M}_{n}(\mathrm{K})\) into a set E, where K is a commutative field, such that, for every triple of matrices \(X\) , \(Y\) , \(Z\) in \(\mathbf{M}_{n}(\mathrm{K})\) ,

\[
\mathbf {F} (X Y Z) = \mathbf {F} (X Z Y).
\]

We except the case \(\mathbf{M}_{2}(\mathbf{F}_{2})\) ; show then that there exists a mapping \(\Phi\) of K into E such that \(\mathrm{F}(X)=\Phi(\det(X))\) . (Using the Corollary to Proposition 17 (no. 9), show first that, for \(U\in\mathbf{SL}_{n}(\mathrm{K})\) , \(\mathrm{F}(X)=\mathrm{F}(XU)\) ; deduce that, in \(\mathbf{GL}_{n}(\mathrm{K})\) , \(\mathrm{F}(X)\) depends only on \(\det(X)\) . On the other hand, if, for \(r\leqslant n-1\) , we write

\[
X _ {r} = \left( \begin{array}{c c} I _ {r} & 0 \\ 0 & 0 \end{array} \right)
\]

show that there exists a matrix \(Y\) of rank \(r\) such that \(X_{r-1} = YX_r\) and \(Y = X_rY\) and deduce that \(F(X)\) has the same value for all matrices of rank \(< n\) .

\begin{enumerate}
\def\labelenumi{\arabic{enumi}.}
\setcounter{enumi}{22}
\tightlist
\item
  In every A-algebra E, if \(x_{1}, \ldots, x_{n}\) are elements of E, we write
\end{enumerate}

\[
\left[ x _ {1} x _ {2} \dots x _ {n} \right] = \sum_ {\sigma \in \mathfrak {S} _ {n}} \varepsilon_ {\sigma} x _ {\sigma (1)} x _ {\sigma (2)} \dots x _ {\sigma (n)}.
\]

Show that if E admits a finite generating system over A, there exists an integer N such that \([x_{1}x_{2}\ldots x_{N}] = 0\) for all elements \(x_{j} (1 \leqslant j \leqslant N)\) of E.

\begin{enumerate}
\def\labelenumi{\arabic{enumi}.}
\setcounter{enumi}{23}
\tightlist
\item
  Let \(X_{i}\) ( \(1 \leqslant i \leqslant m\) ) be square matrices of the same order \(n\) over the commutative ring A. Show that if \(m\) is even then (Exercise 23)
\end{enumerate}

\[
\operatorname{Tr} \left[ X _ {1} X _ {2} \dots X _ {m} \right] = 0
\]

and if \(m\) is odd

\[
\operatorname{Tr} \left[ X _ {1} X _ {2} \dots X _ {m} \right] = m. \operatorname{Tr} \left(X _ {m} \left[ X _ {1} X _ {2} \dots X _ {m - 1} \right]\right).
\]

(If \(\lambda\) is the cyclic permutation \((1,2,\ldots,m)\) , C the cyclic subgroup of \(\mathfrak{S}_m\) generated by \(\lambda\) and H the subgroup of \(\mathfrak{S}_m\) consisting of the permutations leaving \(m\) invariant, note that every permutation in \(\mathfrak{S}_m\) can be written uniquely as \(\sigma\tau\) with \(\sigma\in\mathrm{H}\) and \(\tau\in\mathrm{C}\) ; then use the fact that the trace of a product of matrices is invariant under cyclic permutation of the factors.)

¶ 25. Let A be a commutative ring containing the field Q of rational numbers and let X be a square matrix of even order \(2n \geqslant 2\) over A. Let L be the subset \(\{1, 2, \ldots, n + 1\}\) of \(\{1, 2, \ldots, 2n\}\) and L' its complement. Show the following formula (``Redei's identity''):

\[
(*) \quad \det (X _ {\mathrm{L}, \mathrm{L}}) \det (X _ {\mathrm{L} ^ {\prime}, \mathrm{L} ^ {\prime}}) = \sum_ {\mathrm{H}, \mathrm{K}} (- 1) \frac {r - 1}{n \binom {n - 1} {r}} \det (X _ {\mathrm{H}, \mathrm{K}}) \det (X _ {\mathrm{H} ^ {\prime}, \mathrm{K} ^ {\prime}})
\]

where H runs through the set of subsets of n elements of \(\{1,2,\ldots,2n\}\) containing l and K the set of subsets of n elements of L and \(r=\operatorname{Card}(H\cap L')\) . (Evaluate the coefficient of a term \(x_{\sigma(1),1}x_{\sigma(2),2}\ldots x_{\sigma(2n),2n}\) in the right-hand side of (*) for an arbitrary permutation \(\sigma\in S_{2n}\) , noting that this coefficient can only be \(\neq0\) if \(H=\sigma(K)\) ; show that it can only be \(\neq0\) if \(\sigma(L)=L\) .)

\begin{enumerate}
\def\labelenumi{\arabic{enumi}.}
\setcounter{enumi}{25}
\tightlist
\item
  \begin{enumerate}
  \def\labelenumii{(\alph{enumii})}
  \tightlist
  \item
    In the notation of Proposition 19 (no. 10), suppose that there exist two A{[}X{]}-linear mappings \(g_1, g_2\) of M{[}X{]} into M'{[}X{]} such that \(\psi' \circ g_2 = g_1 \circ \psi\) ; show that there then exists an A{[}X{]}-linear mapping \(g\) of \(M_u\) into \(M_{u'}\) such that \(\phi' \circ g_1 = \phi' \circ \bar{g}\) . Further, if \(g_1\) and \(g_2\) are A{[}X{]}-isomorphisms of M{[}X{]} onto M'{[}X{]}, \(g\) is an A{[}X{]}-isomorphism of \(M_u\) onto \(M_{u'}\) .
  \end{enumerate}
\end{enumerate}

\begin{enumerate}
\def\labelenumi{(\alph{enumi})}
\setcounter{enumi}{1}
\tightlist
\item
  In the notation of no. 10, the endomorphisms \(u, u'\) are called equivalent if there exist two isomorphisms \(f_1, f_2\) of M onto M such that \(u' \circ f_2 = f_1 \circ u\) ; they are called similar if there exists an isomorphism \(f\) of M onto M' such that \(u' \circ f = f \circ u\) . Deduce from (a) that, for \(u\) and \(u'\) to be similar, it is necessary and sufficient that the endomorphisms \(X - \bar{u}\) and \(X - \bar{u}'\) (of the A{[}X{]}-modules M{[}X{]} and M'{[}X{]} respectively) be equivalent.
\end{enumerate}

